\section{Recommendations, Transfer and Limitations}\label{sec:recs}

\textbf{Recommendations.} (1)~Before any release, run the pilot the proposal planned, with real students' own sentences and choices, accuracy reported per group, and a record of how often nothing fits (\plNoFitPct\% of exam sentences). (2)~Treat the learned ranker as unproven: in the pilot, compare it with equal weights on real choices and keep it only if it wins; on our exam it did not. (3)~Keep asking back: it covered most parsing errors at about \cPerSentence{} questions or offers per sentence. (4)~Improve splitting and number reading.

\textbf{Transfer to engineering scheduling.} The same structure fits wherever requirements are written in words and a schedule must satisfy hard constraints. A planner's note such as ``line 2 is down on Tuesday morning; keep changeovers low this week'' would become machine availability (hard) and changeover cost (soft); a solver with a time limit would replace exact enumeration, as a factory has too many schedules to list~\cite{pinedo2022}; and the ranker would learn from the plans planners actually chose, ordering feasible schedules by throughput, idle time and overtime. Misreads cost more there, so confirmation matters more. Text-to-optimisation work~\cite{ramamonjison2023nl4opt,ahmaditeshnizi2024optimus} could seed such a parser; none of this is tested.

\textbf{Limitations.}\label{sec:limits} No real student took part: all sentences, personas, labels and exams were AI-written, so the results show the method works on AI-written students, not that it helps real ones. The two exam annotators agreed on every parsing label, but they were the same model reading the same rules, and the exam writer removed \sExamDropped{} sentences the rules could not settle, so real sentences will be harder; Claude also wrote part of the training data and every persona, so the exams may share its style. The ranking exam has only \sRankExamPersonas{} personas, so its intervals are wide. The data covers one semester and \dUnits{} units, walking times are estimated, and none of the proposal's usability targets were measured.
